%% reviewer_briefing_preread.tex
%% 評価者向け 事前配布資料（Phase 3b スクリーニング）
%% 内容の正は docs/reviewer_briefing_preread.md。こちらは配布用の組版。
%% ビルド: lualatex --interaction=nonstopmode reviewer_briefing_preread.tex
\documentclass[a4paper]{ltjsarticle}

\usepackage{luatexja-fontspec}
\usepackage[margin=20mm,top=18mm,bottom=20mm]{geometry}
\usepackage{booktabs}
\usepackage{tabularx}
\usepackage{array}
\usepackage{xcolor}
\usepackage{enumitem}
\usepackage[hidelinks]{hyperref}

%% --- 配色 ---
\definecolor{accent}{HTML}{1F4E79}
\definecolor{soft}{HTML}{F0F4F8}
\definecolor{warn}{HTML}{FDF1E7}
\definecolor{rule}{HTML}{C8D4E0}

%% --- 見出し ---
\usepackage{titlesec}
\titleformat{\section}{\normalfont\Large\bfseries\color{accent}}{\thesection.}{0.6em}{}
\titleformat{\subsection}{\normalfont\large\bfseries\color{accent}}{\thesubsection}{0.6em}{}
\titlespacing*{\section}{0pt}{1.6\baselineskip}{0.6\baselineskip}
\titlespacing*{\subsection}{0pt}{1.1\baselineskip}{0.4\baselineskip}

%% --- 強調ボックス ---
\newcommand{\keybox}[1]{%
  \par\medskip
  \noindent\colorbox{soft}{%
    \begin{minipage}{\dimexpr\linewidth-2\fboxsep\relax}
      \vspace{2pt}\centering\bfseries\large #1\vspace{4pt}
    \end{minipage}}%
  \par\medskip}

\newcommand{\notebox}[1]{%
  \par\medskip
  \noindent\colorbox{warn}{%
    \begin{minipage}{\dimexpr\linewidth-2\fboxsep\relax}
      \vspace{4pt}#1\vspace{4pt}
    \end{minipage}}%
  \par\medskip}

\newcolumntype{L}[1]{>{\raggedright\arraybackslash}p{#1}}

\setlist[itemize]{leftmargin=1.4em,itemsep=2pt,topsep=3pt}
\renewcommand{\arraystretch}{1.25}

\begin{document}

\begin{center}
  {\LARGE\bfseries\color{accent} 事前お読みいただく資料}\\[4pt]
  {\large Phase 3b\quad Title/Abstract スクリーニング}\\[10pt]
  {\small 説明会: 2026年8月20日（30分）\quad|\quad 作成: 2026年8月19日（8月20日 改訂）}
\end{center}

\vspace{4pt}
\noindent\textcolor{rule}{\rule{\linewidth}{0.8pt}}
\vspace{6pt}

\noindent
いただいたご質問のうち、\textbf{文章で答えられるものを先にお答えします。}
当日は「ファイルの見方」と「実際の論文での判断すり合わせ」に時間を使わせてください。

\noindent
詳細版は \texttt{reviewer\_briefing.md} にあります（長いので、当日の話が分かりにくかった箇所だけ後から引いてください）。

%% ============================================================
\section{何を調べる研究か}

\subsection*{いただいた理解}

\begin{quote}
VR空間においてユーザが自身の身体の大きさをどのように知覚するか
\end{quote}

\noindent
\textbf{ほぼその通りです。} そのうえで、本研究の主張はもう一段先にあります。

\medskip
\noindent
VRでアバタを巨大化・縮小化したとき、脳は2通りの解釈ができます。

\begin{center}
\begin{tabularx}{\linewidth}{L{7em} L{13em} X}
\toprule
 & 解釈 & 例 \\
\midrule
\textbf{自己スケール} & \textbf{自分が}大きく（小さく）なった & 「巨人になった気がする」 \\
\textbf{環境スケール} & \textbf{世界が}小さく（大きく）なった & 「ミニチュアの街を見ている」 \\
\bottomrule
\end{tabularx}
\end{center}

\noindent
この2つは\textbf{主観的にはほとんど区別できません。} そして——ここが本研究の出発点ですが——
\textbf{多くの先行研究が測っているのは後者（環境スケール）のほう}です。
「自己の身体を操作して、外界の見え方を測る」という構造的なズレが分野全体にあります。

\medskip
\noindent
さらに、操作の手段が\textbf{ほぼ視覚に限られている}という偏りもあります。そこで本レビューは、

\begin{enumerate}[leftmargin=1.8em]
  \item この\textbf{視覚偏重と測定のズレを、網羅的な文献調査で定量的に示し}、
  \item 聴覚・触覚といった\textbf{非視覚的手がかりが果たしうる役割を、多感覚統合の理論から導く}
\end{enumerate}

\noindent
ことを目指しています。

\keybox{「VRで自分の体の大きさをどう感じるか」は、これまでほぼ視覚だけで、\\
しかも「自分」ではなく「世界」を測る形で研究されてきた。その偏りを可視化する。}

\subsection*{具体的にどういう論文か（実例）}

著者が事前に「これは確実に対象」と特定している論文が17件あります。そのうち3件を挙げます。

\begin{center}
\begin{tabularx}{\linewidth}{L{9.5em} L{10em} X}
\toprule
論文 & 掲載先 & 何をした研究か \\
\midrule
Leyrer et al. 2011
& \textbf{ACM SIGGRAPH SAP}\newline （HCI系国際会議）
& 眼の高さとアバタの有無を分離して操作し、\textbf{自己中心距離の推定}への影響を測った。\textbf{操作するのは自分の身体、測るのは外界}という、本研究が問題にするズレの典型例 \\
\addlinespace
Serino et al. 2020\newline \textit{Gulliver's virtual travels}
& \textbf{Cognitive Processing}\newline （認知科学系ジャーナル）
& 極端な体サイズへの能動的身体化。縮小では身体イメージが有意に減少し、拡大では変化なし。\textbf{数少ない「自己」を測った研究} \\
\addlinespace
Okada et al. 2025
& \textbf{IEEE VRW}\newline （HCI系国際会議\newline ・\mbox{\textbf{2ページ}}）
& セルフタッチの触覚フィードバックと体サイズの交互作用。\textbf{視覚以外}を扱う数少ない例 \\
\bottomrule
\end{tabularx}
\end{center}

\noindent
ここから読み取っていただきたいのは3点です。

\begin{enumerate}[leftmargin=1.8em]
  \item \textbf{心理・認知系の雑誌と、HCI系の国際会議の両方が対象}です。「HCIの論文だけ」ではありません。
  \item \textbf{2ページの短い論文も対象}です。ページ数は除外理由になりません（実験があれば対象）。
  \item \textbf{視覚だけを扱った研究こそが主要な調査対象}です。「多感覚を扱っていないから物足りない」ということはありません。\textbf{それらが「視覚偏重」という主張の分母になります。}
\end{enumerate}

%% ============================================================
\section{なぜこの手順が必要か}

\subsection*{システマティックレビューという形式の要求}

通常の文献紹介は「著者が良いと思った論文を選んで紹介する」ものです。
これに対しシステマティックレビューは、\textbf{選び方を先に決めて機械的に適用し、
選定の過程を第三者が再現・監査できる形で公開する}ことを要求します。
投稿先の ACM Computing Surveys と、報告基準である PRISMA 2020 がこれを求めています。

\medskip
\noindent
だから「著者が読んで決める」では足りず、\textbf{複数人が独立に判定した記録}が要ります。

\subsection*{なぜ複数人なのか}

1人でスクリーニングすると、\textbf{関連文献の約13\%を見落とす}という無作為化比較試験の実測があります（感度86.6\%。2名なら97.5\%）。

\begin{quote}\small
Gartlehner et al. (2020) \textit{J Clin Epidemiol} 121:20--28. DOI: 10.1016/j.jclinepi.2020.01.005
\end{quote}

\noindent
ただし全1,052件を2名で見ると工数が単純に倍になります。そこで採っているのが
\textbf{liberal accelerated} という方式です。

\keybox{1人が Include にすれば通す。Exclude するには2人必要。}

\noindent
この段階のミスは\textbf{非対称}だからです。

\begin{itemize}
  \item \textbf{誤って Exclude} → その論文は二度と検討されない（\textbf{回復不能}）
  \item \textbf{誤って Include} → 次の全文評価で落とせるだけ（\textbf{手間が増えるだけ}）
\end{itemize}

\noindent
除外の方向にだけ2名を要求することで、工数を抑えつつ見落としを防いでいます。

%% ============================================================
\section{164件と1,052件の関係}

いただいたご質問は

\begin{quote}
164件の後に著者が1,052件を判定する、という流れが十分に理解できておらず
\end{quote}

\noindent
というものでした。\textbf{この順番の理解は正確ではありません。} 正しくはこうです。

\keybox{164件は、1,052件の「一部」です。前後関係ではなく、包含関係です。}

\begin{center}
\begin{minipage}{0.95\linewidth}
\ttfamily\small
1,052件 ── 判定対象のすべて\\
\hspace*{2.0em}│\\
\hspace*{2.0em}├─ \textbf{164件（15.6\%）「校正セット」}\\
\hspace*{2.0em}│\hspace*{2.2em}→ 著者・お二人の \textbf{3名全員} が、同じ164件を独立に判定する\\
\hspace*{2.0em}│\hspace*{2.2em}→ 評価者間一致度（$\kappa$）はこの164件だけで算出する\\
\hspace*{2.0em}│\\
\hspace*{2.0em}└─ \textbf{888件}\\
\hspace*{4.6em}→ stage 1 は著者ひとりが判定する\\
\hspace*{4.6em}→ そのうち \textbf{著者が Exclude / Unsure にしたものだけ} を\\
\hspace*{6.6em}stage 2 で第2評価者が確認する（最大 449件 / 439件）
\end{minipage}
\end{center}

\medskip
\noindent
\textbf{時系列で書くと:}

\begin{center}
\begin{tabularx}{\linewidth}{L{5.5em} L{5em} X L{9em}}
\toprule
 & 誰が & 何件 & いつ \\
\midrule
\textbf{stage 1} & 著者 & \textbf{1,052件すべて}（164件を含む） & いま並行して進行 \\
\textbf{stage 1} & お二人 & \textbf{164件}（校正セット） & いま（〜8/26） \\
\textbf{stage 2} & お二人 & 著者が Exclude / Unsure にした分だけ & 著者の stage 1 完了後 \\
\bottomrule
\end{tabularx}
\end{center}

\noindent
つまり——

\begin{itemize}
  \item \textbf{著者が1,052件を判定するのは「164件の後」ではなく、同時並行です。}
  \item \textbf{お二人の164件は、その1,052件の中から抜き出した部分集合です。} 別物ではありません。
  \item \textbf{164件が終わったら第2ラウンド（stage 2）があります。} 件数は著者の判定が終わるまで確定しません。
\end{itemize}

\subsection*{なぜ164件だけ3人で見るのか}

\textbf{評価者間一致度（$\kappa$）を計算するためです。}
$\kappa$ は「2人の判定がどれだけ一致したか」を偶然の一致を割り引いて測る指標で、
\textbf{同じ対象を複数人が判定していないと計算できません。}

\medskip
\noindent
stage 2 の除外プールだけで計算しようとすると、そこに入っているのは定義上すべて
「著者が Exclude / Unsure にしたもの」なので、\textbf{著者側の判定に分散がなく $\kappa$ が構造的に常に0}に
なってしまいます。そのため、3名全員が同じ集合を判定する部分が別に必要で、それが校正セット164件です。

\notebox{\textbf{この164件も、判定基準は他の888件とまったく同じです。}
「テストだから厳しく」といった調整は不要で、普段どおり判定してください。}

%% ============================================================
\section{8月20日に説明すること}

ご依頼いただいた3点を、この順で扱います。

\subsection{ファイルの説明（Excel と CSV の違いなど）}

\keybox{お二人が触るのは \texttt{sheet\_<お名前>.xlsx} の1つだけです。}

\begin{center}
\begin{tabularx}{\linewidth}{L{13em} X L{7.5em}}
\toprule
ファイル & 何か & 触るか \\
\midrule
\textbf{\texttt{sheet\_<お名前>.xlsx}} & \textbf{記入用の作業ファイル。} ドロップダウン・色分け・入力チェックが入っている & \textbf{これだけ記入する} \\
\texttt{sheet\_<お名前>.csv} & 同じ内容の機械可読な元データ。集計スクリプトが読む正本 & 触らない \\
\texttt{EXAMPLE\_記入見本.xlsx} & 記入例5件。全員に同じものを配布 & 参照のみ \\
\texttt{assignment.csv} & 誰がどれを担当するかの割当表 & 触らない \\
\bottomrule
\end{tabularx}
\end{center}

\noindent
\textbf{なぜ両方あるのか。} CSV は文字だけの単純な形式で、集計プログラムが確実に読めます。
一方 CSV にはドロップダウンも色分けもセルのロックも作れません。
そこで\textbf{「記入は Excel、集計は CSV」}と役割を分けています。
xlsx をご返送いただければ、こちらで CSV に戻して集計します。

\subsection{記入の粒度}

\textbf{「判定」シートの3列だけです。}

\begin{center}
\begin{tabularx}{\linewidth}{L{6em} L{11em} X}
\toprule
列 & 必須か & 粒度の目安 \\
\midrule
\textbf{判定} & \textbf{必須} & \texttt{Include} / \texttt{Exclude} / \texttt{Unsure} から選ぶだけ \\
\addlinespace
\textbf{除外理由} & \textbf{Exclude のときのみ必須} & 9択のドロップダウンから\textbf{1つ}選ぶ。複数該当するときは\textbf{最も明白なもの1つ}を選び、残りはメモへ \\
\addlinespace
\textbf{メモ} & 任意\newline （除外理由が「その他」のときのみ必須） & \textbf{1行で十分。} 後日の協議で「なぜそう判断したか」が復元できる程度。長文は不要 \\
\bottomrule
\end{tabularx}
\end{center}

\noindent
\textbf{メモの粒度の具体例:}

\begin{itemize}
  \item 良い例（Unsure に付ける）:「HMD か明記なし。要旨だけでは判断不能」
  \item 良い例（Exclude・その他）:「同一著者の同一実験の別報告に見えるが確証なし」
  \item 不要:「この研究は〜という背景のもと〜を目的として\ldots」（要約は書かなくて構いません）
\end{itemize}

\noindent
\textbf{判定の原則を1つだけ覚えていただければ十分です。}

\keybox{除外できると確信できないものは残す（Include に倒す）。迷ったら Unsure。}

\noindent
無理に二択にしないでください。\texttt{Unsure} は後日の協議に回る正しい経路です。

\subsection{シートの文言について、1点だけ補足させてください}

除外理由の選択肢に \texttt{S: 原著論文でない} があり、シートの説明には
「総説・サーベイ・抄録のみ・\textbf{ポスター}・基調講演・書籍など」と書かれています。
この「ポスター」の解釈を明確にします。

\par\medskip
\noindent
\begin{minipage}{\linewidth}
\keybox{\texttt{S: 原著論文でない} は\\「\textbf{実験を報告していない}文献種別」を指します。}
\begin{itemize}
  \item \textbf{該当する:} 総説・サーベイ・基調講演・書籍、\textbf{実験の記述がない}抄録
  \item \textbf{該当しない: 実験を報告しているポスター/ショートペーパー → \texttt{Include}}
  \item \textbf{ページ数は除外理由になりません。判断するのは内容です。}
\end{itemize}
\end{minipage}

\medskip
\noindent
第1節で挙げた Okada et al.\ 2025 は2ページの論文ですが、\textbf{対象です。}
著者が「これは確実に対象」と特定している17件のうち\textbf{4件が2ページの会議論文}で、
そのうち1件は本レビューで\textbf{唯一の聴覚の事例}です。ここを落とすと主張の根拠が痩せます。

\notebox{シート側の文言は、配布済みのものを\textbf{そのまま残してあります。}
判定の途中で基準そのものを書き換えると、変更前後で適用された基準が違うことになり、
評価者間一致度（$\kappa$）が意味を失うためです。\textbf{解釈はこの補足が優先します。}}

\subsection{実際の論文でのすり合わせ}

ご要望のあった「1〜2件の論文を例に、該当するかどうかと議論の観点をすり合わせたい」を
当日のメインにします。\textbf{ひとつだけお願いがあります。}

\keybox{例に使うのは、お二人が担当する164件の「外」の論文に\\ 限らせてください。}

\noindent
担当分の論文を事前に一緒に検討してしまうと、その論文について\textbf{独立した判定ではなくなり}、
$\kappa$（評価者間一致度）が意味を失います。stage 2 で配られる可能性のある論文も同様です。

\medskip
\noindent
そこで当日は、\textbf{判定対象にまったく含まれていない論文}を例に使います。
すり合わせたい観点は同じように扱えますし、独立性も保てます。具体的には——

\begin{itemize}
  \item \texttt{EXAMPLE\_記入見本.xlsx} の5件（Include / Exclude 3種 / Unsure）
  \item 上記17件のうち、判定対象1,052件にまったく含まれていない\textbf{9件}\\（第1節に挙げた3件はすべてこれに該当します）
\end{itemize}

\notebox{\textbf{判定中に個別の論文で迷われた場合も、ご相談ではなく \texttt{Unsure} にしてください。}
それが協議に回る正しい経路です（基準そのものへのご質問は、いつでも歓迎です）。}

%% ============================================================
\section{締切}

\keybox{2026年8月26日（火）}

\noindent
記入済みの \texttt{sheet\_<お名前>.xlsx} をご返送ください。集計と $\kappa$ の算出はこちらで行います。

\medskip
\noindent
164件のうち\textbf{35件は要旨がなくタイトルだけ}です（DOIから外部補完しても取得できなかったもの）。
判断できなければ \texttt{Unsure} にしてください。件数は論文の Threats to Validity に明記します。

\medskip
\noindent
要旨のある129件は\textbf{中央値1,466字}あります。1件1〜2分として3〜5時間程度が目安ですが、
迷ったものを \texttt{Unsure} に倒していただければその分短くなります。日をまたいで分割して構いません。

\notebox{\textbf{間に合わない見込みが立った時点でご連絡ください。}
未記入のまま提出されるより、期日を調整するほうが安全です（未記入があると $\kappa$ が算出できません）。}

\end{document}
